\section{Experimental design}
\label{sec:study}

The pre-execution protocol and \texttt{analysis/configs/main.json} freeze the design below. Section~\ref{sec:results} reports achieved coverage, reference checks, predictive diagnostics and measured costs; raw checkpoints and the evidence ledger preserve their provenance.

\paragraph{Data and fitted models.}
We use Statlog German Credit, ACSIncome from California's 2018 one-year ACS PUMS, and OpenML Diabetes (data ID 37, version 1), with 20, 10 and 8 original-feature players, respectively \citep{hofmann1994,uscensus2018pums,ding2021retiring,openml37}. German Credit explains the probability of bad credit (original class 2); ACSIncome explains income above \$50,000; Diabetes explains the positive label. The ACS task uses folktables 0.0.12 and its working-adult filter: age greater than 16, income greater than \$100, positive weekly work hours and person weight at least one. Survey weights do not weight the benchmark. This is a person-row prediction experiment, not a population-representative ACS analysis. Diabetes retains the loader's declared treatment of zero plasma glucose, blood pressure, skin thickness, insulin and body mass index as missing; pregnancy counts of zero remain valid.

Each dataset uses a fixed stratified 80/20 training/test split; ACS training is capped at 20,000 by stratified subsampling. Preprocessing is fitted only on training rows: median imputation and standardization for numeric fields, most-frequent imputation for categorical fields, and one-hot encoding for logistic regression or ordinal encoding for tree models. The model families are logistic regression, a 100-tree random forest and 100-iteration histogram gradient boosting. They provide contrasting CPU-feasible predictors; this is neither a model-selection benchmark nor a survey of all tabular architectures. Training seed 20261103 is fixed. Test accuracy, ROC AUC and log loss are retained for every fitted model, separately from attribution quality.

\paragraph{Inputs, targets and estimators.}
A uniform draw of 32 training rows defines each fixed background. A uniform held-out panel is shared across model types within a dataset; the first two selected inputs are reserved for pilots and excluded from the 60 main inputs. No selection uses explanation outcomes. Each input/estimator/budget has 30 independently seeded runs, with model, background, feature groups and probability output fixed. Seeds derive from base 314159 and dataset/model/input/estimator/budget/run identifiers. The design compares grouped Kernel SHAP (SHAP 0.50.0, identity link, \texttt{l1\_reg=0}) with antithetic permutation sampling, which averages a random feature order and its reverse. It does not evaluate every available Shapley estimator.

Requested coalition-work settings are 128/512/2048 for German Credit, 64/256/512 for ACSIncome, and 32/64/128 for Diabetes. These dimension-specific settings are intended to span finite approximation budgets while avoiding needless exhaustion of smaller games. Permutation runs use an even number of paths and cache repeated coalitions within a run. Kernel SHAP records the effective varying-player count, sampled coalitions and design-rank diagnostics. Nominal settings need not execute equal model-row work: background expansion, cached coalitions, endpoints, integer paths and exhaustive cells all affect cost. Comparisons therefore require actual prediction rows and measured CPU/elapsed time alongside the nominal setting. Reference and pilot costs belong in total study accounting, separately from one-run cost.

The accepted reference strategy enumerates every coalition of every selected empirical game, using the same model-output adapter and grouped features as the approximations. Batched prediction bounds memory. Exact enumeration is a target-accuracy control; an efficiency check alone cannot validate its coordinates. Appendix~\ref{app:implementation-checks} reports completed numerical checks and an unsuccessful exploratory control, alongside the completed main reference audit. The completed matrix contains nine fitted dataset/model cells, 540 reference games and 97,200 stochastic calls. The saved-panel audit verifies these achieved denominators. Missing or failed runs cannot be scored as completed negative results, and neither pilot-only cases nor partial-debug aggregates enter the main evidence.

\paragraph{Operational criteria and uncertainty.}
The primary scenario in Eq.~\eqref{eq:adequacy} fixes $k=3$, $\delta=.01$ probability units, $\tau=.20$ and $\eta=2/3$. It asks for bounded vector error and agreement on at least two signed salient features. It is an explicitly selected numerical tolerance, without a clinical or lending interpretation. Sensitivities include error cutoffs $.10/.20/.40$, complete signed top-three recovery, a $10^{-6}$ probability-unit tie/sign tolerance, normalization floors $.001/.01/.05$ and weak-signal cases. Primary exact ties use original-feature index order. Reference-defined strata distinguish $\|\phi^\star\|_1<.01$ and a normalized largest-versus-second-largest magnitude gap at most $.02$; such strata describe conditions rather than new independent experiments.

Pointwise 95\% descriptive intervals use 2,000 instance-cluster percentile bootstrap resamples (seed 8127) of per-input run means. Paired contrasts preserve input identity. The inference is conditional on the fitted model, split and background; it does not average over training or background resampling. Intervals are not simultaneous across the grid. Boundary-rate handling and the limits of finite-panel certification are stated in Section~\ref{sec:assessment}. A separate post-main background-size/resampling arm changes the target and remains outside the fixed-background repetition estimand (Appendix~\ref{app:background}).
